% ---
% id: reductions/rem:claimscope
% kind: rem
% title: "The claims that read a declared number"
% summary: >-
%   The squashing, transfer, calculus and core instance results read no
%   declared number; a submodule's instance reads declared data decided
%   exactly, the licence exactly Hermitian operators and declared constants.
%   Where the normalizer acts, every statement is read on the declared model.
% prominence: standard
% status: accepted
% uses: [thm:emit, def:emitnorm, rem:emitclaim, def:reduction-instance]
% mentions: [cor:identity-instance, cor:flag-instance, thm:instrument-instance, cor:register-instance, thm:licence, thm:chain, req:emission-record]
% assumes: [A-CONST]
% refs: []
% provenance: newly-derived
% lean:
% ---

\begin{remark}\label{rem:claimscope}
\emph{The \tref{squashing relation} and its soundness, the factorization and
transfer theorems, the calculus of the pre-order, the branching rule, the
flags-enough corollary, register-conditioned squashing and the obstruction
read no declared number} and assert nothing about the emitted artefact: each
is a statement about an arbitrary squashing relation, an arbitrary
instrument, an arbitrary register-conditioned mixture, or an explicit
counterexample. \emph{The instance theorems of the core} --- the identity,
flag, branching and register instances
(Corollaries~\ref{cor:identity-instance}, \ref{cor:flag-instance}
and~\ref{cor:register-instance}, Theorem~\ref{thm:instrument-instance}) ---
are statements about arbitrary admissible data and arbitrary branch instances
and read no declared number either. \emph{A submodule's instance} builds its
datum out of declared blocks and declared classical data; its hypotheses are
conditions on a POVM family and on declared scalars and matrices, and this
module's refusal rules decide every one of those conditions \emph{exactly},
so it applies to the declared data verbatim. \emph{The attack-set licence}
(Theorem~\ref{thm:licence}) needs its emitted constraint operators exactly
Hermitian and its constants declared: the first holds by construction and is
decided by the refusal rule that says so, the second by
Assumption~\ref{ass:constdecl}. \emph{The chain theorem}
(Theorem~\ref{thm:chain}) asserts $\textsf{status}=\textsf{exact}$ of the
artefact, which is true of the numbers emitted
(Theorem~\ref{thm:emit}\,(iii)). \emph{The emission theorem}
(Theorem~\ref{thm:emit}) is the claim about all of the above, stated once over
the reduction instance (Definition~\ref{def:reduction-instance}).

\emph{One conditioning survives and it is the only one}: where the emission
normalizer acts, each of those statements is read on the declared model of
Definition~\ref{def:emitnorm} --- on $\Bsys$, the declared model
$\Mdec=M[A']$; in the data sectors, the declared datum $A'$ --- and not on
the handed datum, and Remark~\ref{rem:emitclaim} is what that means.
\emph{Every residual is recorded}: the displacement and the rounding
directions are (Requirement~\ref{req:emission-record}), and the consumer's
key length is Theorem~\ref{thm:emit}\,(v) --- not established, and no refusal
rule can establish it.
\end{remark}
